\chapter{Differences from the 1992 Manual}\label{app:differences}
\sloppy

The standard for both implementations is the 1992 manual,
\emph{Dylan: an object-oriented dynamic language} (\textsl{DIRM}). Where
an implementation departs from it --- by necessity, by a decision where
the manual is silent or contradicts itself, or by adding something ---
this appendix says so. It collects in one place what the other chapters
mention as they go. Each entry says which implementation it concerns:
\emph{both}, the \emph{interactive} Common Lisp implementation, or the
\emph{compiler} to C.

\section{Different behavior, in both implementations}
\label{sec:diff-both}

\begin{description}
\item[\code{!key}, \code{!rest} and \code{!next}.]
  \index{"!key@\code{"!key}}
  These are read as \code{\#key}, \code{\#rest} and \code{\#next}, as
  Thomas read them, because Thomas's programs --- the book's examples
  among them --- use them. So they cannot be used as variable names.
  Every other name the manual allows is allowed.
  \code{((method (a !key (b 2)) (list a b)) 1 b: 5)} \yields{}
  \code{(1 5)}.

\item[Modules.] \index{modules}
  The 1992 manual describes modules but leaves the way they are declared
  to the programming environment. Both implementations declare them with
  \code{define-module}, as design note~32 describes, written in prefix
  form, and name a file's module with a \code{module:} header, as design
  note~33 describes. Chapter~\ref{ch:modules} has the details.

\item[Backquote.] \index{backquote}
  Chapter 3 of the manual defers backquote to ``a separate section of
  this manual'' that the manual does not contain. Both implementations
  follow the usual Lisp rules for \code{`}, \code{,} and \code{,@},
  nesting included: \code{`(1 ,(+ 1 1) ,@(list 3 4))} \yields{}
  \code{(1 2 3 4)}.

\item[\code{map-into}.] \index{map-into@\code{map-into}}
  The function is passed only the elements of the source collections, as
  chapter 12's prose and design note~27 say. The chapter's one example,
  \code{(map-into x + '(1 2 3 4))} with \code{x} holding
  \code{(100 100 200 200)}, shows \code{(101 102 203 204)}, which would
  require passing the target's elements too; here it answers
  \code{(1 2 3 4)}.

\item[Irrational functions.]
  \code{sqrt}, \code{exp}, \code{log}, the trigonometric functions and
  their kin answer floats for exact arguments, as in Common Lisp, which
  the manual cites for them, and as the manual's examples show:
  \code{(sqrt 25)} \yields{} \code{5.0}, \code{(exp 0)} \yields{}
  \code{1.0}. Rationals are otherwise exact: \code{(/ 6 4)} \yields{}
  \code{3/2}.

\item[Method specificity among unrelated classes.]
  \index{specificity}
  The manual orders applicable methods by comparing their specializers,
  but leaves open the order of two classes neither of which is a
  subclass of the other. Both implementations follow CLOS: the class
  that comes first in the \emph{argument's} class precedence list is
  more specific.
\begin{dylan}
(define-class <a> (<object>))
(define-class <b> (<object>))
(define-class <c> (<a> <b>))
(define-method which ((x <a>)) 'a)
(define-method which ((x <b>)) 'b)
(which (make <c>))                 ; => a
\end{dylan}
\end{description}

\section{Additions in both implementations}\label{sec:diff-additions}

\begin{description}
\item[Printing.] \code{display}, \code{newline}, \code{print},
  \code{write-line}, \code{write} and \code{pp}, and the generic
  function \code{print-object}, to which programs add methods to control
  how their objects print (Chapter~\ref{ch:printing}). The manual uses
  \code{print} in its examples without defining it.
\item[Streams.] \code{<stream>}, \code{*standard-output*},
  \code{*standard-error*}, \code{<string-stream>} with
  \code{stream-contents}, and \code{format}; every printing function
  takes \code{stream:}. The manual uses streams and \code{format} but
  does not define them. There is no input.
\end{description}

\section{The interactive implementation}\label{sec:diff-interactive}

\begin{description}
\item[Lisp access.] \code{lisp-variable} and \code{lisp-procedure}
  reach Common Lisp values and functions from Dylan. Lisp's \code{NIL}
  arrives as the empty list, not \code{\#f}
  (Chapter~\ref{ch:lisp}).
\item[\code{\#t} and \code{\#f}.] Their class, as \code{object-class}
  answers it, is \code{<object>}; the manual names no class for them.
  (The compiler differs; see below.)
\item[No extensions of the compiler.] The interactive implementation has
  none of the compiler's macros and threads.
\end{description}

\section{The compiler to C}\label{sec:diff-compiler}

\subsection{What a static program cannot do}

The compiler compiles whole, static programs, whose classes and generic
functions are fixed when they are compiled
(Section~\ref{sec:dylanc-static}). The manual allows a running program
to change them; a compiled program cannot:

\begin{itemize}
\item \code{add-method}, \code{remove-method}, \code{add-slot},
  \code{remove-slot} and \code{freeze-methods} are compile errors,
  wherever they appear.
\item Each specializer of a method made with \code{define-method}, and
  each superclass of a class, must be known when compiling: a class, a
  constant variable holding one, \code{(singleton \meta{x})}, or
  \code{(make <singleton> object: \meta{x})} of a constant. A
  \code{define-slot} owner may be any singleton, but a class owner must
  be known when compiling.
\item A class cannot be defined twice.
\item \code{(make <class> ...)} is an error when it runs. A
  \code{<generic-function>} made by \code{make} has no methods, and none
  can be added.
\item There is no evaluation or loading of code at run time.
\end{itemize}

\subsection{Numbers}

\begin{itemize}
\item \textbf{Integers are 63 bits}, from $-2^{62}$ to $2^{62}-1$. There
  are no bignums. A literal out of range is a compile error; an
  arithmetic result out of range signals ``integer overflow (integers are
  63 bits)'' rather than giving a wrong answer. The numerator and
  denominator of a ratio have the same limit. The interactive
  implementation has bignums: \code{(* 4611686018427387904 2)}
  \yields{} \code{9223372036854775808} there, while in a compiled program
  even the literal is an error.
\item \textbf{No complex numbers.} \code{<complex>} exists, abstract, as
  the superclass of \code{<real>} (so \code{(instance? 3 <complex>)}
  \yields{} \code{\#t}), but it has no instances. A complex literal is a
  compile error, and a function whose result would be complex signals an
  error: \code{(sqrt -4.0)} signals ``sqrt: the result would be complex''
  where the interactive implementation answers \code{\#c(0.0 2.0)}.
\item \code{<single-float>} is a C \code{float}; \code{<double-float>}
  and \code{<extended-float>} are both a C \code{double}. Floats print
  as the shortest decimal that reads back: \code{21.0}, and \code{1.5s0}
  for a single float.
\end{itemize}

\subsection{Conditions without a debugger}

There is no interactive debugger. An unhandled serious condition prints
\code{Dylan error:} and its message on the standard error and ends the
program with status~1 --- or, in a thread other than the main one, ends
that thread (Chapter~\ref{ch:extensions}). An unhandled warning prints
\code{Dylan warning:} and its message and continues. \code{break}
prints its message and returns \code{\#f}. \code{restart-query} and
\code{return-query} do nothing.

\subsection{Representation}

\code{\#t} and \code{\#f} are instances of an internal class,
\code{<boolean>}, which \code{object-class} answers but which the
\code{dylan} module does not export (the manual names no class for
them). \code{seal} works as the manual says, and \code{make-read-only}
answers its argument, a class, unchanged, since a compiled program's
classes are already fixed.

\subsection{Additions}

All of these are extensions; the 1992 manual has none of them.

\begin{itemize}
\item \textbf{Macros}: \code{define-macro}, \code{define-for-macros},
  and \code{gensym} at compile time; \cmd{dylanc -E} to see expansions
  (Chapter~\ref{ch:extensions}).
\item \textbf{Threads}: \code{make-thread}, \code{join-thread},
  \code{current-thread}, \code{thread-yield}; \code{<lock>},
  \code{<recursive-lock>}, \code{acquire-lock}, \code{release-lock} and
  \code{with-lock}; \code{<condition-variable>} with
  \code{wait-condition-variable}, \code{notify-condition-variable} and
  \code{notify-all-condition-variable}; \code{<atomic-integer>} with
  \code{atomic-value}, \code{atomic-add!} and
  \code{atomic-compare-and-swap!}; \code{define-thread-variable} and
  \code{dynamic-bind}; \code{<thread-error>} with
  \code{thread-error-thread} and \code{thread-error-condition}
  (Chapter~\ref{ch:extensions}).
\item \textbf{Timing}: \code{current-milliseconds}.
\end{itemize}

\section{Where the two implementations disagree}\label{sec:diff-between}

Besides the differences above, the implementations disagree where the
manual leaves a choice open or where the host shows through. A portable
program depends on none of these; each is described where it arises.

\begin{itemize}
\item The operator of a call is evaluated first by the compiler to C
  and last by the interactive implementation; arguments are evaluated
  left to right in both (Section~\ref{sec:syntax-calls}).
\item A method whose parameter list has \code{\#key} with no keywords
  accepts no keywords in the compiler to C and any keyword in the
  interactive implementation (Section~\ref{sec:functions-keyword-checking}).
\item A handler that returns from the error \code{check-type} signals
  leaves it unhandled in the compiler to C; in the interactive
  implementation \code{check-type} answers what the handler returned
  (Section~\ref{sec:conditions-signalling}).
\item Equal string literals may be one object in the compiler to C;
  in the interactive implementation each is its own
  (Section~\ref{sec:collections-strings}).
\item \code{min} and \code{max} of equal numbers of different classes,
  and functions whose true result is complex or infinite, answer
  differently (Sections~\ref{sec:numbers-misc}
  and~\ref{sec:numbers-transcendental}).
\item Assigning a read-only module variable is a compile error in the
  compiler to C and an error when it runs in the interactive
  implementation, which also needs modules declared before use
  (Sections~\ref{sec:modules-readonly} and~\ref{sec:modules-lisp}).
\item Error messages are worded differently
  (Section~\ref{sec:conditions-runtime-errors}).
\end{itemize}

\fussy
